\section{Kontrola}

\begin{frame}{Před exportem do PDF}
  \begin{columns}[T,totalwidth=\textwidth]
    \begin{column}{.48\textwidth}
      \textbf{Obsah}
      \begin{enumerate}
        \item Každý slide má jeden účel.
        \item Nadpis sděluje hlavní závěr.
        \item Text je stručný a bez překlepů.
        \item Obrázky něco vysvětlují.
      \end{enumerate}
    \end{column}
    \begin{column}{.48\textwidth}
      \textbf{Technická kontrola}
      \begin{enumerate}
        \item Překlad končí bez chyb.
        \item Nic nepřetéká mimo slide.
        \item Patička a číslování sedí.
        \item PDF funguje i na jiném počítači.
      \end{enumerate}
    \end{column}
  \end{columns}
  \vspace{4mm}
  \begin{Warning}
    Neupravuj barvy a rozměry přímo ve slidech. Patří do souborů v \texttt{theme/}.
  \end{Warning}
\end{frame}

\begin{frame}{Kdy rozšířit samotnou šablonu}
  \begin{enumerate}
    \item Nejdřív zkus běžný slide, dva sloupce nebo velký obrázek.
    \item Opakovaný prvek přidej do \texttt{vsak-components.sty}.
    \item Novou barvu definuj v \texttt{vsak-colors.sty}.
    \item Titulku, sekce a patičku měň v \texttt{beamerthemeVSAK.sty}.
  \end{enumerate}
  \vspace{4mm}
  \begin{Rule}
    Obsah používá pojmenované komponenty. Jejich vzhled se řídí centrálně, takže redesign nevyžaduje přepisovat všechny slidy.
  \end{Rule}
\end{frame}
